\documentclass[11pt,a4paper]{article}
\usepackage[utf8]{inputenc}
\usepackage[T1]{fontenc}
\usepackage[portuguese]{babel}
\usepackage{xcolor}
\usepackage{hyperref}
\usepackage{geometry}
\usepackage{listings}
\usepackage{tcolorbox}
\geometry{margin=2.5cm}
\definecolor{primary}{RGB}{41,128,185}
\definecolor{secondary}{RGB}{44,62,80}
\definecolor{codebg}{RGB}{240,240,240}
\lstset{backgroundcolor=\color{codebg},basicstyle=\ttfamily\small,breaklines=true,frame=single}
\title{\color{secondary}\Huge\textbf{Solução: Exercício 2.6 - Monitoramento básico}}
\author{\color{primary}DevOps com Kubernetes -- 2026}
\date{}
\begin{document}
\maketitle


\textbf{Guia de Solução:}

\subsection*{\color{secondary}Passo a Passo}

1. Verifique métricas rápidas com \texttt{kubectl top pods} e \texttt{kubectl top nodes} quando houver Metrics Server.
2. Para Prometheus/Grafana, prefira namespace \texttt{monitoring} e instalação via Helm.
3. Observe reinícios com \texttt{kubectl get pods} e detalhes com \texttt{kubectl describe pod}.
4. Registre quais métricas indicam saúde da aplicação: CPU, memória, READY, restarts e latência.

\end{document}
